\uestcsymbols

% 每项依次填写符号、说明和符号在正文中首次出现的页码。
% 推荐使用\pageref引用正文中的标签，使页码能够自动更新。
% 多个符号在正文同一位置首次出现时，可以引用同一个标签。
% 以下符号只用于演示常见写法；请删除论文中未实际使用的条目。
% 内容较多时，可以手动拆分表格或自行改用支持跨页的表格环境。
\begingroup
\renewcommand{\arraystretch}{1.3}
\setlength{\tabcolsep}{9pt}
\noindent
\begin{tabularx}{\linewidth}{%
    >{\normalsize}l
    >{\normalsize\raggedright\arraybackslash}X
    >{\normalsize}c}
    \textbf{符号} & \textbf{说明} & \textbf{页码} \\
    $\pi$ & 圆周率 & \pageref{sym:math-symbols} \\
    $\theta$ & 旋转角 & \pageref{sym:math-symbols} \\
    $\mathbf{x}$ & 输入向量 & \pageref{sym:math-symbols} \\
    $\mathbf{y}$ & 输出向量 & \pageref{sym:math-symbols} \\
    $\mathbf{A}$ & 变换矩阵 & \pageref{sym:math-symbols} \\
\end{tabularx}
\par
\endgroup
